\documentclass[11pt,twoside,leqno]{article}
\usepackage[T1]{fontenc}
\usepackage{amsmath,amssymb,amsthm,mathtools}
\usepackage{txfonts}
\usepackage[a4paper,textwidth=15cm,textheight=21cm,centering]{geometry}
\usepackage{microtype,needspace,etoolbox}
\usepackage[colorlinks=true,hyperfootnotes=false,linkcolor=blue,citecolor=blue,urlcolor=blue]{hyperref}
\allowdisplaybreaks
\pagestyle{myheadings}
\markboth{\normalfont\small\scshape Sharp vanishing order}
         {\normalfont\small\scshape Biharmonic equations with bounded potentials}
\setcounter{tocdepth}{1}
% The manuscript is anonymous; no author or affiliation is inferred.
\makeatletter
\renewcommand{\@maketitle}{%
  \newpage\null\vskip 1em
  \begin{center}
    {\LARGE\bfseries \@title\par}
  \end{center}
  \vskip 1em}
\makeatother
\setlength{\parindent}{13pt}
\setlength{\emergencystretch}{2em}
\numberwithin{equation}{section}
\newtheorem{theorem}{Theorem}[section]
\newtheorem{proposition}[theorem]{Proposition}
\newtheorem{lemma}[theorem]{Lemma}
\newtheorem{corollary}[theorem]{Corollary}
\theoremstyle{definition}
\newtheorem{definition}[theorem]{Definition}
\theoremstyle{remark}
\newtheorem{remark}[theorem]{Remark}
\BeforeBeginEnvironment{theorem}{\Needspace{5\baselineskip}}
\BeforeBeginEnvironment{proposition}{\Needspace{5\baselineskip}}
\BeforeBeginEnvironment{lemma}{\Needspace{5\baselineskip}}
\newcommand{\R}{\mathbb R}
\newcommand{\C}{\mathbb C}
\newcommand{\ord}{\operatorname{ord}}
\newcommand{\supp}{\operatorname{supp}}
\newcommand{\loc}{\mathrm{loc}}
\newcommand{\eps}{\varepsilon}
\newcommand{\Om}{\Omega}
\hypersetup{pdftitle={Sharp vanishing order for biharmonic equations with bounded potentials},pdfauthor={}}

\title{Sharp vanishing order for biharmonic\\
equations with bounded potentials}
\author{}
\date{}

\begin{document}
\maketitle
\begingroup
\renewcommand{\thefootnote}{}
\footnotetext{\textit{2020 Mathematics Subject Classification.}
Primary 35B60; Secondary 35J30, 35J40.\par
\textit{Key words and phrases.} Biharmonic equation, vanishing order,
bounded potential, quantitative unique continuation, clamped replacement.}
\endgroup
\begin{center}
\begin{minipage}{11.5cm}\small
\noindent\textbf{Abstract.}
The exponent $1/3$ in normalized vanishing-order bounds for biharmonic
equations with bounded potentials is optimal in every dimension $n\ge2$,
even when both the solution and the potential are real. For each
sufficiently large even integer $N$, we construct a
bounded real solution of $\Delta^2u_N+V_Nu_N=0$ on the plane with
$\ord_0u_N=N+2$, $\|V_N\|_\infty\le CN^3$, and
$\|u_N\|_{L^\infty(B_1)}=\|u_N\|_{L^\infty(\R^2)}=1$.
The construction joins biharmonic modes across annuli and removes the
residual near the real zeros by clamped replacement on sparse sets.
Signed radial averaging with vanishing second moment preserves the
replacement near every zero away from the origin, without a maximum
principle.
The solution is globally $C^4$, and both the solution and the potential
are smooth away from the origin. The same family also establishes
sharpness in the complex-valued classes.
\end{minipage}
\end{center}
\medskip
\tableofcontents
\clearpage

\section{Introduction and main results}\label{sec:introduction}

We study the dependence of vanishing order on the size of a bounded
potential in the biharmonic equation
\begin{equation}\label{eq:equation}
 \Delta^2u+Vu=0,
 \qquad \Delta=\sum_{j=1}^n\partial_j^2.
\end{equation}
A bound in terms of the potential alone requires a normalization of the
solution. Here the domain and two comparison scales are fixed, and the
ratio of the outer and inner supremum norms is controlled independently
of the potential. Under this normalization, we determine the smallest
power of $\|V\|_\infty$ in a uniform bound for the vanishing order.
The exponent is $1/3$ in every dimension $n\ge2$. The examples proving
sharpness have both real solutions and real potentials; hence the same
exponent is optimal in the real, real-potential, and fully complex
classes.

\subsection{Background and related results}

The question belongs to quantitative unique continuation. Strong unique
continuation excludes zeros of infinite order, whereas a quantitative
estimate controls the finite order in terms of the coefficients and an
appropriate normalization. The second-order theory provides two
complementary approaches: singular-weight Carleman inequalities, as in
Jerison and Kenig~\cite{JK}, and frequency-function arguments, developed
by Garofalo and Lin~\cite{GL86,GL87}. Both approaches relate local
vanishing to estimates comparing the size of a solution at different
scales. The normalization is part of such an estimate, rather than a
choice that can be made after the constants have been fixed.

Spectral problems suggest a natural scale for this question. For a
nonzero eigenfunction of $-\Delta_g$ on a fixed smooth, closed,
connected Riemannian manifold, Donnelly and Fefferman~\cite{DF}
bounded the vanishing order by a constant multiple of the square root
of the eigenvalue. Kukavica~\cite{Kuk} studied the corresponding
question for variable potentials. For bounded potentials in the
second-order equation, the quantitative estimates of Bourgain and
Kenig~\cite{BK}, arising in Anderson localization, lead instead to the
power $2/3$. Kenig~\cite[p.~XX--4]{Kenig} formulated the normalized
local problem explicitly and asked whether the power could be reduced
to $1/2$ for real solutions and real potentials.

The distinction between these powers reflects both the coefficient
class and the role of complex values. Meshkov~\cite{Meshkov}
constructed complex solutions and bounded complex potentials with decay
of order $\exp(-c|x|^{4/3})$ at infinity. Bakri~\cite{Bakri}
proved quantitative vanishing estimates on closed manifolds and gave a
complex construction showing the sharpness of the bounded-potential
power $2/3$. Under differentiability assumptions on the potential,
stronger bounds involve its derivatives: see Bakri~\cite{Bakri12}
and the weighted-frequency approach of Zhu~\cite{Zhu16}. Estimates
for singular lower-order coefficients were developed further by Davey
and Zhu~\cite{DZ} and Davey~\cite{Davey}. These results distinguish
bounds controlled only by a potential norm from bounds that also depend
on additional regularity.

For powers of the Laplacian, strong unique continuation with singular
lower-order coefficients was studied by Colombini and
Grammatico~\cite{CG} and Lin~\cite{Lin07}. Lin, Nagayasu and
Wang~\cite{LNW} obtained quantitative estimates through three-sphere
and doubling inequalities. Higher-order frequency methods were
considered in~\cite{Zhu16}, and Zhu~\cite{Zhu} established
quantitative estimates with coefficients in Lebesgue spaces. In the
biharmonic case, the iterated Carleman inequality of Huang, Wang and
Zheng~\cite[Lemma~2.1]{HWZ} has a factor $\lambda^6$ on the
left-hand side. It absorbs a bounded zeroth-order term once
$\lambda\ge C(1+\|V\|_\infty^{1/3})$, which identifies the
bounded-potential exponent relevant here. Their separate sharpness
example concerns a different fourth-order constant-coefficient
operator, not $\Delta^2$.

Teng, Wang and Zhu~\cite[Theorem~2]{TWZ} obtain the same
bounded-potential power $1/3$ by weighted frequency functions, with
normalization on two fixed balls. We prove the upper bound for the
precise class used here in Section~\ref{sec:upper}, directly from
the Carleman inequality in~\cite{HWZ}; thus the upper estimate is
independent of the construction. It applies to complex potentials and
requires no derivatives of $V$. This last point matters when comparing
with stronger-regularity results. Under their normalization, Liu, Tian
and Yang~\cite[Theorem~1.1]{LTY} assume $V\in W^{1,\infty}$
and obtain a bound involving both $\|V\|_\infty^{1/4}$ and
$\|\nabla V\|_\infty$. Such a bound does not yield a
$1/4$-power estimate in the $L^\infty$ norm alone. Likewise, the
biharmonic spectral problems treated by Lin and Zhu~\cite{LZ} have
specified boundary conditions and a spectral parameter; they are not
the local problem with an arbitrary bounded potential considered here.

Reality presents a further issue, separate from regularity. In the
second-order planar problem, Kenig, Silvestre and Wang~\cite{KSW}
proved quantitative Landis estimates under a sign condition on the
potential. Logunov, Malinnikova, Nadirashvili and Nazarov~\cite{LMNN}
removed the sign condition in their planar result on exponential decay,
using specifically two-dimensional methods. In contrast, Frank and
Ivanisvili~\cite{FI} constructed real counterexamples with
$4/3$-power exponential decay in dimensions at least three. Feng, Gan,
Jiang and Lin~\cite[Theorem~1.1]{FGJL} constructed real Dirichlet
solutions in a fixed three-dimensional ball with vanishing order $k$
and potential norm at most $Ck^{3/2}$. These results concern different
normalizations and, in the global case, additional control at infinity.
They motivate the real-valued construction but are not substitutes for
a sharpness example in the fixed two-scale class below.

The obstruction to obtaining a real example by taking real parts is
already visible in the equation. If $u=u_1+iu_2$ and $V=V_1+iV_2$,
then $\Delta^2u_1+V_1u_1=V_2u_2$. Thus $u_1$ need not solve a
homogeneous scalar equation with a bounded real potential. Taking real
parts also introduces zeros at which the quotient defining the potential
may become singular. We instead build a real approximate solution and
remove its residual near the entire small-value set. Sparse-set
replacement and variable-radius averaging follow the general strategy
of~\cite[Sections~6--9]{FI} and~\cite[Section~4]{FGJL}, but the
fourth-order correction requires a different inverse estimate and a
different averaging kernel. Positivity is not available as a general
principle for clamped biharmonic problems; see~\cite{GGS}. Our inverse
is obtained from a planar $H^2$ estimate, and the averaging kernel is
chosen to reproduce biharmonic functions exactly.

\subsection{The normalized classes}

Write $B_r^n(p)$ for an open Euclidean ball and $B_r^n=B_r^n(0)$.
All supremum norms are essential suprema. For a locally bounded function,
set
\begin{equation}\label{eq:order}
 \ord_pu=\sup\{\gamma\ge0:
 \|u\|_{L^\infty(B_r(p))}=O(r^\gamma)\text{ as }r\downarrow0\}.
\end{equation}
The implied constant may depend on the individual function and on
$\gamma$. The zero need not be isolated.

\begin{definition}\label{def:classes}
For $n\ge2$, $C_0\ge1$ and $M\ge1$, let
$\mathcal A_n^{\C}(C_0,M)$ consist of the pairs satisfying
\begin{equation}\label{eq:class}
 \begin{gathered}
 \Delta^2u+Vu=0\quad\text{in }\mathcal D'(B_{10}^n),\\
 V\in L^\infty(B_{10}^n;\C),\qquad \|V\|_\infty\le M,\\
 u\in H^4_{\loc}(B_{10}^n;\C)\cap L^\infty(B_{10}^n),\\
 \|u\|_{L^\infty(B_{10}^n)}\le C_0,\qquad
 \|u\|_{L^\infty(B_1^n)}\ge1.
 \end{gathered}
\end{equation}
The subclass with real $u,V$ is denoted by $\mathcal A_n^{\R}$;
that with real $V$ and possibly complex $u$ is denoted by
$\mathcal A_n^{\R,\C}$. No boundary condition or sign restriction
on a real potential is imposed.
For each class, let $\alpha_n$ with the corresponding superscript
be the infimum of all $\alpha\ge0$ such that, for every $C_0\ge1$,
there is $C=C(n,C_0,\alpha)$ for which
\begin{equation}\label{eq:admissible}
 \ord_pu\le C(1+M^\alpha),\qquad p\in B_{1/2}^n,
\end{equation}
holds for every $M\ge1$ and every pair in that class.
\end{definition}

The two-scale condition cannot be omitted. The harmonic polynomials
$10^{-N}\operatorname{Re}(x_1+ix_2)^N$ have outer supremum one and
arbitrarily large vanishing order with zero potential, but do not
satisfy the inner normalization.

\subsection{Main results}

The main construction is planar and entirely real. The common upper
bound then identifies the optimal exponent in each normalized class.

\begin{theorem}[Real examples]\label{thm:real}
There are absolute constants $C>0$ and $N_0$ such that, for every even
integer $N\ge N_0$, there are real functions
\[
 u_N\in C^4(\R^2)\cap L^\infty(\R^2),\qquad
 V_N\in L^\infty(\R^2)
\]
satisfying, pointwise and distributionally,
\begin{equation}\label{eq:family}
 \Delta^2u_N+V_Nu_N=0,\qquad
 \|V_N\|_\infty\le CN^3,\qquad \ord_0u_N=N+2.
\end{equation}
Moreover,
\begin{equation}\label{eq:normalization}
 \|u_N\|_{L^\infty(\R^2)}
 =\|u_N\|_{L^\infty(B_{10}^2)}
 =\|u_N\|_{L^\infty(B_1^2)}=1.
\end{equation}
Both functions are smooth on $\R^2\setminus\{0\}$. With
$z=x_1+ix_2$, there is a nonzero real $d_N$ such that
\begin{equation}\label{eq:jet}
 D^\beta\bigl(u_N-d_N|z|^2\operatorname{Re}(z^N)\bigr)
 =O_N(|z|^{N+6-|\beta|}),\qquad |\beta|\le4.
\end{equation}
\end{theorem}

\begin{proposition}[The common upper bound]\label{prop:upper}
For every $n\ge2$ and $C_0\ge1$, there are $c,C>0$, depending only
on $n,C_0$, such that every pair in $\mathcal A_n^{\C}(C_0,M)$ satisfies
\begin{equation}\label{eq:lower-bound}
 \|u\|_{L^\infty(B_r^n(p))}
 \ge cr^{C(1+M^{1/3})},\qquad p\in B_1^n,\quad0<r<1/4.
\end{equation}
In particular, $\ord_pu\le C(1+M^{1/3})$ on $B_{1/2}^n$.
\end{proposition}

\begin{corollary}\label{cor:sharp}
For every $n\ge2$,
\[
 \alpha_n^{\R}=\alpha_n^{\R,\C}=\alpha_n^{\C}=\frac13,
\]
and the estimate is valid at the endpoint $\alpha=1/3$ in all three
classes. The same real family, extended constantly in the last $n-2$
coordinates, proves sharpness in each class.
\end{corollary}

The theorem supplies the lower obstruction in the smallest of the
three classes, while Proposition~\ref{prop:upper} gives the upper
estimate in the largest. Constant extension preserves the equation,
the normalization, and the vanishing order. No separate complex-valued
construction is therefore needed. The assertion is sharpness of the
power in a uniform estimate, not of its multiplicative constant.

These examples rule out smaller exponents, but do not contradict strong
unique continuation: each zero has finite order, and the potential norm
grows with that order. Neither global smoothness nor compact support
of $V_N$ is asserted. The result concerns $\Delta^2+V$ with bounded
zeroth-order potentials and fixed two-scale normalization; it does not
assert the same exponent for arbitrary fourth-order elliptic operators.
The distinction from prescribed decay at infinity is discussed in
Section~\ref{sec:remarks}.

\subsection{Proof mechanism}

The construction begins with annular transitions from frequency
$\ell$ to $\ell+2k$, with residual bounded by $C\ell^2k^2r^{-4}$
relative to a positive envelope. The two radial branches
$r^{-p}\cos(p\theta)$ and $r^{2-p}\cos(p\theta)$ allow the profile
to remain exactly biharmonic throughout the region of possible
amplitude cancellation. A finite frequency schedule ensures
$\ell k\le N^{3/2}$ and keeps $\sum k^{-1}$ bounded, so that the
transitions fit in a fixed annulus. Kelvin inversion then produces the
prescribed polynomial leading term at the origin.

To control the quotient near the real zeros, we correct the resulting
function $f_N$ on
$D_N=\{|f_N|<\tau E_N\}$. A uniform density estimate makes $D_N$ sparse on balls of radius
$\mu_N^{-1}$. The clamped inverse gives a single correction satisfying
\[
 |w_N|\le C\nu^2\tau E_N,
 \qquad
 |w_N|\le C\nu^2N^3\mu_N^{-4}E_N.
\]
The first estimate keeps the smoothed function away from zero outside
the smaller small-value region contained in $D_N$; the second preserves
the leading polynomial and controls the fourth-order residual after
smoothing. The averaging kernel has vanishing second moment, so it
reproduces the biharmonic replacement near every zero away from the
origin. On the remaining region, the lower bound relative to the
envelope controls the quotient potential.

There are two distinct cancellations in this procedure. The radial
branches eliminate the residual where their amplitudes can cancel;
the zero second moment eliminates the Laplacian term in the radial
mean-value formula for a biharmonic function. The latter kernel is
necessarily signed. Its purpose is exact reproduction, not positivity.
The clamped correction is solved on the sparse open set and extended
by zero as an $H^2$ function. No equation for the fourth derivatives of
that extension across the boundary is assumed.

The proof follows these dependencies. Section~\ref{sec:preliminaries}
records the radial and Kelvin identities.
Section~\ref{sec:approximate} constructs the approximate solution and
its positive envelope, and Section~\ref{sec:density} proves the
uniform small-value density estimate. The inverse and averaging lemmas
are established in Section~\ref{sec:tools} and combined in
Section~\ref{sec:completion} to produce the bounded real potential.
Section~\ref{sec:upper} proves the upper bound and the passage to
higher dimensions independently of the construction.

\section{Preliminaries}\label{sec:preliminaries}

Throughout the construction, $\Om=\R^2\setminus\{0\}$.
We use $D^jf$ for the tensor of all ordered Cartesian derivatives of
order $j$, and
$\|D^2f\|_2^2=\sum_{i,j}\|\partial_{ij}f\|_2^2$.
For an arbitrary open set $D\subset\R^2$,
$H^2_0(D)$ is the closure of $C_c^\infty(D)$ in $H^2(D)$.
Zero extension maps this space into $H^2(\R^2)$, as follows by
approximation. This notation imposes no regularity on $\partial D$.
On a smooth boundary, the clamped traces are $w=\partial_\nu w=0$.

Unless otherwise indicated, the constants $C,C_j$ in the construction
depend only on the fixed smooth profiles and on $j$, and are uniform
in the frequencies and radii. All auxiliary parameters are fixed before
$N$ is allowed to vary; the constants in \eqref{eq:jet} may depend on
$N$.

\begin{lemma}[Radial identities]\label{lem:radial}
With $t=\log r$ and $D=\partial_t$,
\begin{equation}\label{eq:polar}
 \Delta^2=e^{-4t}\bigl((\partial_t-2)^2+\partial_\theta^2\bigr)
                         (\partial_t^2+\partial_\theta^2).
\end{equation}
For an integer $m\ge1$,
\begin{equation}\label{eq:conjugation}
 \Delta^2[r^{-m}e^{im\theta}b(t)]
 =r^{-m-4}e^{im\theta}P_m(D)b,
\end{equation}
where
\begin{align*}
 P_m(D)&=D(D-2)(D-2m)(D-2m-2)\\
 &=D^4-(4m+4)D^3+(4m^2+12m+4)D^2-(8m^2+8m)D.
\end{align*}
Consequently, for $1\le k\le m$,
\begin{equation}\label{eq:radial-cost}
 |P_m(k\partial_x)b|
 \le Cm^2k^2\max_{1\le j\le4}\|b^{(j)}\|_\infty.
\end{equation}
Also, for real $s$,
\begin{equation}\label{eq:powers}
 \Delta^2(r^se^{im\theta})
 =(s^2-m^2)((s-2)^2-m^2)r^{s-4}e^{im\theta}.
\end{equation}
\end{lemma}
\begin{proof}
The identity $\Delta=e^{-2t}(\partial_t^2+\partial_\theta^2)$
gives \eqref{eq:polar}. Conjugating by $e^{-mt}$ replaces
$\partial_t$ by $D-m$, which yields \eqref{eq:conjugation}.
The derivative terms are bounded using
$k^4+mk^3+m^2k^2+m^2k\le Cm^2k^2$.
Applying the two factors in \eqref{eq:polar} to $e^{st+im\theta}$
gives \eqref{eq:powers}.
\end{proof}

\begin{lemma}[Kelvin transformation]\label{lem:kelvin}
For a smooth function $v$ on a planar annulus, set
$Kv(r,\theta)=r^2v(1/r,\theta)$. Then
\begin{equation}\label{eq:kelvin}
 \Delta^2(Kv)(r,\theta)=r^{-6}(\Delta^2v)(1/r,\theta).
\end{equation}
\end{lemma}
\begin{proof}
Write $g(s,\theta)=v(e^s,\theta)$ and $s=-t$.
Applying \eqref{eq:polar} to $e^{2t}g(-t,\theta)$ gives $e^{-2t}$
times the commuting factors
$\partial_s^2+\partial_\theta^2$ and
$(\partial_s-2)^2+\partial_\theta^2$ acting on $g$.
The corresponding expression for $(\Delta^2v)(e^{-t},\theta)$
has the prefactor $e^{4t}$. Comparing the two expressions proves the
identity.
\end{proof}

\section{Construction of the approximate solution}\label{sec:approximate}

We construct annular profiles with controlled residuals, then choose
a frequency schedule that keeps all transitions in a fixed annulus.

\subsection{Annular transitions}

Fix $\delta=10^{-3}$ and real smooth cutoffs with values in $[0,1]$:
\[
\begin{array}{c|cc}
 &\text{equals }0&\text{equals }1\\
 \chi_2(x)&x\le-\delta&x\ge\delta\\
 \chi_1(x)&x\ge3/4&x\le1/4\\
 \eta(x)&x\le1&x\ge7/4.
\end{array}
\]
For integers $\ell\ge16$, $1\le k\le\ell$, and $\rho>0$, put
\begin{equation}\label{eq:annular-parameters}
 p=\ell+2k,\quad r_c=\rho e^{1/k},\quad r_+=\rho e^{3/k},\quad
 x=k\log(r/r_c)\in[-1,2].
\end{equation}
Set
\[
 L_k(x)=20k(e^{2x/k}-1),\qquad
 H_k(x)=(1-\eta(x))L_k(x)+\eta(x)L_k(2),\qquad
 Q_k(x)=e^{-2x}L_k(x).
\]
Define
\begin{equation}\label{eq:profile}
 F_{\ell,k,\rho}
 =\chi_1(x)r^{-\ell}\cos(\ell\theta)
 -\chi_2(x)r_c^{2k}r^{-p}H_k(x)\cos(p\theta).
\end{equation}
Write
\[
 A=r^{-\ell},\quad b_1=\chi_1,\quad
 b_2=\chi_2e^{-2x}H_k,\quad Z=b_1-b_2e^{2ik\theta},\quad a=k/\ell.
\]
Choose $\omega\in C_c^\infty((0.002,0.2))$, with
$0\le\omega\le1$ and $\omega=1$ on $[0.005,0.15]$, and set
\begin{equation}\label{eq:envelope}
 E_{\ell,k,\rho}=A(|Z|^2+a^2\omega(x)^2)^{1/2}.
\end{equation}
Thus $F_{\ell,k,\rho}=A\operatorname{Re}(e^{i\ell\theta}Z)$.
The complex amplitude is used only as notation; both the profile and
the envelope are real.

\begin{lemma}[Profiles and envelopes]\label{lem:profile}
The envelope is positive and smooth. Uniformly in $\ell,k,\rho$,
\begin{align}
 |F_{\ell,k,\rho}|&\le E_{\ell,k,\rho}\le C_A r^{-\ell},
                                                     \label{eq:profile-size}\\
 |D^jF_{\ell,k,\rho}|+|D^jE_{\ell,k,\rho}|
 &\le C_j(\ell/r)^jE_{\ell,k,\rho},\qquad j\ge1,
                                                     \label{eq:profile-jets}\\
 |\Delta^2F_{\ell,k,\rho}|
 &\le C\ell^2k^2r^{-4}E_{\ell,k,\rho}.                \label{eq:profile-source}
\end{align}
The profile is biharmonic throughout $\delta<x<1/4$. On collars of
the inner and outer boundaries, respectively, the profile and envelope
are
\[
 (r^{-\ell}\cos\ell\theta,r^{-\ell}),\qquad
 (-r_c^{2k}L_k(2)r^{-p}\cos p\theta,r_c^{2k}L_k(2)r^{-p}).
\]
The ratio of the outer to the inner envelope values is at most
$80e^{-27/4}<1$.
\end{lemma}
\begin{proof}
The derivatives of $L_k,H_k,Q_k$ of each fixed order are uniformly
bounded on $[-1,2]$. For $j\ge1$, this follows from
$L_k^{(j)}=20\,2^jk^{1-j}e^{2x/k}$; for order zero, use
$|L_k(x)|\le40|x|e^{2|x|}$.
The identity
\begin{equation}\label{eq:exact-branches}
 r_c^{2k}r^{-p}L_k(x)
 =20kr_c^{2k}(r_c^{-2}r^{2-p}-r^{-p})
\end{equation}
together with \eqref{eq:powers} shows that the profile is exactly
biharmonic in the crossing region, where $\chi_1=\chi_2=1$ and $\eta=0$.

Amplitude cancellation does not occur where the profile cutoffs vary:
\[
 |Q_k(x)|\le40\delta e^{4\delta}<1/4\quad(|x|\le\delta),
 \qquad Q_k(x)\ge40xe^{-2x}>5\quad(1/4\le x\le3/4).
\]
For the second inequality, the only interior critical point of
$xe^{-2x}$ is a maximum, and both endpoint values exceed $1/8$.
On $1\le x\le2$, we also have $H_k\ge L_k(1)\ge40$.
Hence $|Z|\ge c>0$ wherever a profile cutoff varies. Applying
\eqref{eq:radial-cost} there at frequencies $\ell$ and $p\le3\ell$
bounds the residual by
$C\ell^2k^2r^{-4}A\le C\ell^2k^2r^{-4}E$.
The profile is exactly biharmonic in the remaining regions.

A zero of $Z$ requires $Q_k=1$ and $e^{2ik\theta}=1$ in
$\delta<x<1/4$. The representation
\begin{equation}\label{eq:Q-integral}
 Q_k(x)=40x\int_0^1e^{-2x(1-t/k)}\,dt
\end{equation}
gives $40xe^{-2x}\le Q_k(x)\le40x$ for $x>0$. Hence
\begin{equation}\label{eq:crossing-band}
 1/2\le Q_k(x)\le2,\quad 0<x<1/4
 \quad\Longrightarrow\quad 1/80\le x<1/10.
\end{equation}
Indeed, $x\ge1/10$ would give $Q_k\ge4e^{-1/5}>2$.
Moreover,
\begin{equation}\label{eq:Qprime}
 Q_k'(x)=40e^{-2x}\bigl(k-(k-1)e^{2x/k}\bigr),\qquad
 16\le Q_k'(x)\le40\quad(0.005\le x\le0.15).
\end{equation}
For the lower bound, use
$k-(k-1)e^{2x/k}\ge1-2xe^{2x}$ and
$e^{-2x}\ge1-2x$; the upper bound follows from
$k-(k-1)e^{2x/k}\le1$. The same formula shows $Q_k'>0$
on $(0,1/4)$, since for $k>1$,
$\log(1-1/k)+2x/k<0$; the case $k=1$ is direct.
Thus $\omega=1$ wherever $|Z|$ can approach zero, which proves that
the envelope is positive.

Where $|Z|$ is bounded below, \eqref{eq:profile-jets} follows by
differentiation. It remains to treat the cancellation band. Write
$E=AB$ with $B=(|Z|^2+a^2)^{1/2}\ge a$.
Logarithmic-polar derivatives of $Z$ of order $h\ge1$ are bounded by
$C_hk^h$. Derivatives of the square root give terms
$C_jk^jB^{1-h}$, $1\le h\le j$, and
\[
 \frac{k^jB^{1-h}}{\ell^jB}=\frac{a^j}{B^h}\le a^{j-h}\le1.
\]
For $F=A\operatorname{Re}(e^{i\ell\theta}Z)$ use $|Z|\le B$
and $k^h\le\ell^ha\le\ell^hB$. Passing to Cartesian derivatives introduces $r^{-j}$ and terms of
lower logarithmic-polar order, proving \eqref{eq:profile-jets}.
The cutoff $\omega$ varies only where $|Z|\ge c$, so it introduces
no further loss.

The size bound and the collar formulas follow from the definitions.
The ratio of the outer to the inner envelope values is
\[
 L_k(2)e^{-4-3\ell/k}
 \le80e^{-4-(3\ell-4)/k}
 \le80e^{-7+4/\ell}\le80e^{-27/4},
\]
where $L_k(2)\le80e^{4/k}$.
\end{proof}

\subsection{Frequency budget and assembly}

The frequency increments must control both the residual and the total
radial width.
Let $0<\eps\le1/2$ and let $L\ge\max\{16,4\eps^{-2}\}$ be even.
For an even $N>L$, start at $\ell_0=L$, $\rho_0=2$ and set
\begin{equation}\label{eq:schedule}
 \begin{split}
 k_j&=\min\{\lfloor\eps\ell_j\rfloor,
             \lfloor N^{3/2}/\ell_j\rfloor,(N-\ell_j)/2\},\\
 \ell_{j+1}&=\ell_j+2k_j,\qquad \rho_{j+1}=\rho_je^{3/k_j},
 \end{split}
\end{equation}
until $\ell_J=N$.

\begin{lemma}\label{lem:budget}
The schedule terminates, and
\begin{equation}\label{eq:budget}
 1\le k_j\le\eps\ell_j,\quad \ell_j^2k_j^2\le N^3,\quad
 \sum_{j<J}k_j^{-1}<3,\quad 2\le\rho_j\le R_*:=2e^9.
\end{equation}
The radius bound is independent of $N,L,\eps$.
\end{lemma}
\begin{proof}
Before termination, every candidate for $k_j$ is a positive integer:
$\eps L\ge2$ and $N^{3/2}/\ell_j\ge\sqrt N$, while the remaining
gap is positive and even. Thus $\ell_j$ reaches $N$ after finitely
many steps. The first two estimates follow from the definition.

A step strictly limited by the remaining gap is terminal and contributes
at most one to $\sum k_j^{-1}$. Among the other steps, those with
$k_j=\lfloor\eps\ell_j\rfloor$ occur first, since this candidate
increases while the other decreases. For these steps,
$k_j\ge\eps\ell_j/2$ and $\ell_{j+1}\ge(1+\eps)\ell_j$, so
\[
 \sum\frac1{k_j}\le\frac{2}{\eps L}
 \sum_{h\ge0}(1+\eps)^{-h}
 =\frac{2(1+\eps)}{\eps^2L}\le\frac34.
\]
For the remaining steps not limited by the final gap,
$k_j\ge N^{3/2}/(2\ell_j)$ and
\[
 \frac1{k_j}=\frac{\ell_{j+1}-\ell_j}{2k_j^2}
 \le\frac2{N^3}\int_{\ell_j}^{\ell_{j+1}}t^2\,dt.
\]
Their sum is at most $2/3$. Hence the total is less than three,
and the radius recurrence gives \eqref{eq:budget}.
\end{proof}

Set $\gamma_0=2^L$ and
$\gamma_{j+1}=-\gamma_jr_{c,j}^{2k_j}L_{k_j}(2)$, where
$r_{c,j}=\rho_je^{1/k_j}$. On the $j$th annulus use
$\gamma_jF_{\ell_j,k_j,\rho_j}$ and $|\gamma_j|E_{\ell_j,k_j,\rho_j}$.
The collars agree with the same harmonic expression on both sides of
each interface, so the assembled functions match to all orders.
Continue by
\[
 \gamma_Jr^{-N}\cos N\theta,\qquad |\gamma_J|r^{-N}
 \quad(r\ge\rho_J),
\]
and denote the assembled pair by $v,\mathcal E$.
The envelope equals one at the initial boundary. The boundary-ratio
estimate therefore gives
\begin{equation}\label{eq:exterior-envelope}
 |v|\le\mathcal E\le C_A\qquad(r\ge2).
\end{equation}
Every $\gamma_j$ is nonzero.

\subsection{Inversion, outer extension, and local frequency}

Kelvin inversion converts the terminal mode into the prescribed
polynomial at the origin. For $0<r\le1/2$, set
\begin{equation}\label{eq:prepared-inversion}
 f_N(r,\theta)=r^2v(1/r,\theta),\qquad
 E_N(r,\theta)=r^2\mathcal E(1/r,\theta).
\end{equation}
In the fixed core $0<r<R_*^{-1}$, the formulas reduce to
\begin{equation}\label{eq:core}
 f_N=c_Nr^{N+2}\cos N\theta,\qquad
 E_N=|c_N|r^{N+2},\qquad c_N=\gamma_J\ne0.
\end{equation}
Choose a fixed smooth $\zeta\in[0,1]$, zero for $r\le5/8$ and one
for $r\ge3/4$. For $r\ge1/2$, extend by
\begin{equation}\label{eq:outer}
 \begin{split}
 A_L(r)&=2^Lr^{L+2}\exp[-(2L+2)\zeta(r)\log(2r)],\\
 f_N(r,\theta)&=A_L(r)\cos L\theta,\qquad E_N(r,\theta)=A_L(r).
 \end{split}
\end{equation}
The definitions agree on a collar of $r=1/2$, and
$A_L(r)=2^{-L-2}r^{-L}$ for $r\ge3/4$.
Writing $\mathcal L_LA=A''+r^{-1}A'-L^2r^{-2}A$, we have
\begin{equation}\label{eq:outer-source}
 \Delta^2f_N=q_Lf_N,\qquad
 q_L=\frac{\mathcal L_L^2A_L}{A_L},\qquad
 |q_L|\le C(L/r)^4,\qquad \supp q_L\subset[5/8,3/4].
\end{equation}
For the bound, differentiation of $\log A_L$ on the fixed transition
interval gives $|\partial_r^j\log A_L|\le C_jLr^{-j}$ for $j\ge1$.
The product rule then yields $|A_L^{(j)}|\le C_j(L/r)^jA_L$ for $j\le4$;
substitution into $\mathcal L_L^2$ proves \eqref{eq:outer-source}.
Outside this interval, the two exact powers are biharmonic. Only the
positive amplitude $A_L$ appears in the denominator; no division by
the cosine factor is needed.

Fix a smooth $\vartheta\in[0,1]$ with $\vartheta=0$ on $x\le-1/2$
and $\vartheta=1$ on $x\ge7/4$. On each exterior annulus define
$m(r)=\ell_j+2k_j\vartheta(x)$, with $m=N$ on the terminal tail. Set
\begin{equation}\label{eq:frequency}
 \mu_N(r)=m(1/r)/r\quad(r\le1/2),\qquad
 \mu_N(r)=L/r\quad(r\ge1/2).
\end{equation}
The collar identities ensure that this defines a smooth positive
function on $\Om$.
In exterior logarithmic radius, $\partial_t^hm=O_h(k^{h+1})$ for
$h\ge1$. In particular,
\[
 \left|\frac{d}{dr}\frac{r}{m(1/r)}\right|
 \le L^{-1}+C(k/\ell)^2.
\]
Thus fixed $\eps_0>0$ and $L_0\ge16$ can be chosen so that, whenever
$\eps\le\eps_0$ and $L\ge L_0$,
\begin{equation}\label{eq:frequency-jets}
 \mu_N\ge L/r,\qquad |\nabla\mu_N^{-1}|\le1/100,\qquad
 |D^j\mu_N|\le C_j\mu_N^{j+1}\quad(j\ge1).
\end{equation}
For higher derivatives, the preceding bounds for $m$ give
$|\partial_t^hm|\le C_hm^{h+1}$; conversion to radial derivatives,
together with $m\ge L$, proves the last estimate. Decrease $\eps_0$
so that every ball of radius $3/\mu_N$ meeting an assembly interface
lies in the matching pure-mode collars, where the angular frequency
in $\mu_N$ is constant. This choice is uniform: the collars have
logarithmic width at least $c/k$, whereas the logarithmic width of
such a ball is at most $C/\ell$. The same conclusion holds after
inversion.

The following estimates will be used in the density argument and the
clamped correction.

\begin{proposition}[Prepared estimates]\label{prop:prepared}
Assume the schedule hypotheses, $\eps\le\eps_0$, $L\ge L_0$, and
$N^3\ge L^4$. With constants independent of $N,L,\eps$,
\begin{align}
 |f_N|&\le E_N,\qquad
 |D^jf_N|+|D^jE_N|\le C_j\mu_N^jE_N\quad(j\ge1),
                                                    \label{eq:prepared-jets}\\
 |\Delta^2f_N|&\le M_NE_N,\qquad M_N=C_*N^3.         \label{eq:global-source}
\end{align}
For $A=E_N$ and $A=E_N\mu_N^{-4}$ there is one fixed $K\ge1$ such that
\begin{equation}\label{eq:comparison}
 K^{-1}\le\frac{\mu_N(y)}{\mu_N(x)}\le K,\qquad
 K^{-1}\le\frac{A(y)}{A(x)}\le K
 \quad\text{if }|y-x|<4/\mu_N(x).
\end{equation}
These balls are relatively compact in $\Om$. For $0<\tau<1$, on
$D_N(\tau)=\{|f_N|<\tau E_N\}\subset\Om$,
\begin{equation}\label{eq:small-source}
 |\Delta^2f_N|\le C(\eps^2+\tau)\mu_N^4E_N.
\end{equation}
\end{proposition}
\begin{proof}
The jet estimates follow from \eqref{eq:profile-jets}, the derivatives
of inversion, and \eqref{eq:outer}. In an inverted annulus,
\eqref{eq:kelvin} and \eqref{eq:profile-source} give
\[
 |\Delta^2f_N|\le C\ell_j^2k_j^2r^{-4}E_N
 \le CR_*^4N^3E_N,\qquad
 |\Delta^2f_N|\le C\eps^2\mu_N^4E_N.
\]
Here $R_*^{-1}\le r\le1/2$ on every transition. The residual vanishes in the core. On the outer transition,
\eqref{eq:outer-source} and $L^4\le N^3$ give the global bound;
the additional inequality $|f_N|<\tau E_N$ gives
\eqref{eq:small-source}.

Since $\mu_N\ge L/r$, the comparison balls avoid the origin. The
Lipschitz bound for $\mu_N^{-1}$ gives
$0.96\le\mu_N(x)/\mu_N(y)\le1.04$ on these balls.
Moreover, $|\nabla\log E_N|\le C\mu_N$ by \eqref{eq:prepared-jets}.
Integration along line segments proves the comparison for $E_N$.
Combining this with the frequency comparison gives the estimate for
the second weight.
\end{proof}

\section{Uniform density of the small-value set}\label{sec:density}

The clamped correction requires the small-value set to occupy a small
fraction of every ball at the local frequency scale. The density
threshold must be chosen uniformly before $\eps,L$ are fixed.
We prove this by compactness of the normalized profiles. The only
degenerate regime is the coalescence of the two angular modes,
which is described by the following lemma.

\begin{lemma}[Coalescing modes]\label{lem:coalescing}
Let $\ell_j\to\infty$, $1\le k_j\le\ell_j$, and
$a_j=k_j/\ell_j\to0$. Suppose
$x_j\in(\delta,1/4)$, $Q_{k_j}(x_j)\to1$ and $e^{i\psi_j}\to1$.
For arbitrary phases $\phi_j$, set
\[
 Z_j=1-Q_{k_j}(x_j)e^{i\psi_j},\qquad
 b_j=(|Z_j|^2+a_j^2)^{1/2}.
\]
After passage to a subsequence, there are $\phi\in\R$, $c\in\C$,
$d\ge0$ and $q\in[16,40]$ with
\[
 e^{i\phi_j}\to e^{i\phi},\quad Z_j/b_j\to c,\quad
 a_j/b_j\to d,\quad Q_{k_j}'(x_j)\to q,\quad |c|^2+d^2=1,
\]
such that, for every fixed $m\ge0$,
\begin{equation}\label{eq:coalescing}
 \frac{1-Q_{k_j}(x_j+a_jy)e^{i(\psi_j+2a_js)}}{b_j}
 \longrightarrow c-d(qy+2is)
 \quad\text{in }C^m_{\loc}(\R^2_{y,s}).
\end{equation}
The real function
\begin{equation}\label{eq:limit-profile}
 G(y,s)=e^{-y}\operatorname{Re}
       [e^{i(\phi+s)}\{c-d(qy+2is)\}]
\end{equation}
has a zero set of planar measure zero.
\end{lemma}
\begin{proof}
By \eqref{eq:crossing-band}, eventually $x_j\in[1/80,1/10)$.
This interval stays a fixed positive distance from the endpoints of
the crossing region and lies where $\omega=1$.
Equation~\eqref{eq:Qprime}, together with compactness of the phases,
gives the stated subsequential limits. Their normalization follows
from $|Z_j/b_j|^2+(a_j/b_j)^2=1$.

For a fixed compact set $K$ and integer $m$, Taylor's formula gives
\begin{align*}
 Q_{k_j}(x_j+a_jy)e^{i(\psi_j+2a_js)}
 ={}&e^{i\psi_j}Q_{k_j}(x_j)\\
 &+a_je^{i\psi_j}
   (Q_{k_j}'(x_j)y+2iQ_{k_j}(x_j)s)+R_j(y,s),
\end{align*}
where
\begin{equation}\label{eq:Taylor-remainder}
 \|R_j\|_{C^m(K)}\le C_{K,m}a_j^2.
\end{equation}
To verify that the remainder bound is uniform, take two derivatives
with respect to $t$ of
$Q_{k_j}(x_j+ta_jy)e^{i(\psi_j+2ta_js)}$. This gives
\[
 a_j^2e^{i(\psi_j+2ta_js)}
 [y^2Q_{k_j}''+4iysQ_{k_j}'-4s^2Q_{k_j}](x_j+ta_jy).
\]
All its fixed-order $y,s$ derivatives are bounded by $C_{K,m}a_j^2$,
since the evaluation points remain in the crossing region and the
derivatives of $Q_k$ are uniformly bounded. The integral remainder proves
\eqref{eq:Taylor-remainder}. Since $b_j\ge a_j$,
$\|R_j/b_j\|_{C^m(K)}\le C_{K,m}a_j\to0$.
This proves \eqref{eq:coalescing}.

If $d=0$, then $|c|=1$, so $G$ is a nonzero sinusoid multiplied by
$e^{-y}$.
If $d>0$, writing $c=c_1+ic_2$ gives
\[
 e^yG=(c_1-dqy)\cos(\phi+s)-(c_2-2ds)\sin(\phi+s).
\]
For each $s$ with $\cos(\phi+s)\ne0$, this has at most one zero in
$y$. The exceptional values of $s$ form a discrete set. Fubini's
theorem proves the assertion.
\end{proof}

\begin{proposition}[Uniform density]\label{prop:density}
For every $\nu>0$ there is $\tau>0$, independent of $N,L,\eps$ in
Proposition~\ref{prop:prepared}, such that
\begin{equation}\label{eq:density}
 |D_N(\tau)\cap B_{2/\mu_N(X)}(X)|
 \le\nu|B_{2/\mu_N(X)}|\qquad(X\in\Om).
\end{equation}
\end{proposition}
\begin{proof}
In the core and the outer extension,
$|f_N|/E_N=|\cos(m\theta)|$ and $\mu_N=m/r$, with $m=N$ or $L$.
A ball of radius $2r/m$ has radial width at most $4r/m$ and angular
width at most $C/m$. Its angular interval meets at most a fixed number of cosine zeros.
The angular measure of the set where $|\cos(m\theta)|<\tau$ is at
most $C\tau/m$ for $\tau<1/2$. Integrating with the area element
$r'\,dr'\,d\theta$ gives a relative measure at most $C\tau$.
Balls meeting an assembly interface lie in matching pure-mode collars
by the choice of $\eps_0$ and satisfy the same estimate.

For any remaining ball with center $X_0$, use
$\xi=\mu_N(X_0)(X-X_0)$ and divide $f_N$ by $E_N(X_0)$.
Equations~\eqref{eq:prepared-jets} and \eqref{eq:comparison} give
uniform $C^h$ bounds on a neighborhood of $\overline{B_2}$, for each
fixed $h$. The normalized profiles are therefore precompact in every fixed
$C^m$ topology. It remains to identify their limits and show that
each has a zero set of measure zero.

We may work with the exterior annular profiles. Normalization cancels
the assembly coefficient up to a sign, which is fixed after passage
to a subsequence. Inversion and its factor $r^2$ contribute a
nondegenerate change of coordinates and a nonvanishing multiplier.
To verify uniformity, let $(r_0,\theta_0)$ be the exterior center,
set $m_0=m(r_0)$, and write
\begin{equation}\label{eq:local-coordinates}
 y=\ell\log(r/r_0),\quad s=\ell(\theta-\theta_0),\quad
 x=x_0+(k/\ell)y.
\end{equation}
The Jacobians into frequency-rescaled Euclidean coordinates are
\begin{equation}\label{eq:Jacobian}
 (m_0/\ell)^2e^{2y/\ell}
 \quad\text{and, after inversion,}\quad
 (m_0/\ell)^2e^{-2y/\ell}.
\end{equation}
Because $1\le m_0/\ell\le3$, these and their reciprocals are bounded
on fixed bounded coordinate sets. For bounded $\ell$ the maps have
smoothly convergent subsequences; for $\ell\to\infty$ they converge,
after a subsequence, to invertible linear maps. The normalized Kelvin multiplier is
$(r_0/r)^2=e^{-2y/\ell}$ and has a nonvanishing limit.
Neither operation changes whether the zero set has measure zero.

The centers satisfy $x_0\in[-1,2]$. Passing to a subsequence, we may
assume that $x_0$ and all angular phases converge. Every ball not
covered by the collar argument remains within a single annular profile.

\emph{Bounded $\ell$.}
The integers $\ell,k$ are fixed on a subsequence. Positivity of the
envelope and compactness of the center parameters bound
$E(r_0,\theta_0)/r_0^{-\ell}$ below. The limit is a radial combination
of two angular waves of distinct frequencies $\ell$ and $\ell+2k$.
The radial coefficients never vanish simultaneously: wherever the
first can vanish, the second is nonzero. Hence the angular
trigonometric polynomial is nontrivial at each fixed radius. Its
zero set has one-dimensional measure zero, and Fubini's theorem gives
the assertion.

\emph{$\ell\to\infty$ and $k/\ell\to a_*>0$.}
Here $k\to\infty$, $L_k\to40x$ smoothly on bounded intervals, and
$H_k$ has the corresponding smooth limit. The center envelope remains
bounded below after division by $r_0^{-\ell}$: wherever cancellation
is possible, $\omega=1$ and $a\to a_*>0$.
In \eqref{eq:local-coordinates} the limit is a nonzero radial factor
times two angular waves of frequencies $1$ and $1+2a_*$. At each radial coordinate, their coefficients are not both zero.
Since the frequencies are distinct, the angular analytic function is
nontrivial. Fubini's theorem again gives a zero set of planar measure
zero.

\emph{$\ell\to\infty$ and $a=k/\ell\to0$.}
Set $\phi=\ell\theta_0\pmod{2\pi}$,
$\psi_0=2k\theta_0\pmod{2\pi}$ and
$b=E(r_0,\theta_0)/r_0^{-\ell}$.
If $b$ stays bounded below, the radial coefficients and the relative
phase become constant in the limit, while $a^2\omega^2$ tends to
zero. The limit is a nonzero multiple of
$e^{-y}\operatorname{Re}(e^{i(\phi+s)}Z_*)$, with $Z_*\ne0$.
If $b\to0$, the separation of the cutoff regions forces the centers
into the exactly biharmonic crossing region, and
\[
 Q_k(x_0)\to1,\quad e^{i\psi_0}\to1,\quad
 b=(|1-Q_k(x_0)e^{i\psi_0}|^2+a^2)^{1/2}.
\]
Lemma~\ref{lem:coalescing} gives the limit \eqref{eq:limit-profile}.
In either case, the zero set has planar measure zero.

Suppose \eqref{eq:density} fails for a fixed $\nu>0$. Choose $\tau_j\downarrow0$ and corresponding balls on which the
estimate fails, and rescale each ball to $B_2$. A subsequence of the
normalized profiles converges uniformly to $g$ with
$|\{g=0\}\cap B_2|=0$. By envelope comparison, each rescaled
small-value set is contained in $\{|g_j|<K\tau_j\}\cap B_2$. If
$\delta_j=K\tau_j+\|g_j-g\|_{L^\infty(B_2)}$, then
\[
 \{|g_j|<K\tau_j\}\cap B_2\subset\{|g|\le\delta_j\}\cap B_2.
\]
For every $t>0$, the upper limit of these measures is at most
$|\{|g|\le t\}\cap B_2|$. Letting $t\downarrow0$ and using
continuity of measure from above gives zero, contrary to the fixed
positive lower bound. The sets remain in the fixed ball $B_2$
throughout the argument.
\end{proof}

\section{Clamped replacement and biharmonic averaging}\label{sec:tools}

The correction uses two tools: a clamped inverse estimate obtained
without positivity, and a signed radial average that reproduces
biharmonic functions on the sampling balls.

\subsection{A planar sparse inverse}

The first lemma combines a local energy estimate with a pointwise
bound for functions supported on sets of small measure.

\begin{lemma}[Local energy and small supports]\label{lem:energy}
Let $D\subset\R^2$ be open and let $w\in H^2_0(D;\R)$ satisfy
\[
 \int_D\Delta w\Delta\phi=\int_DF\phi
 \quad(\phi\in H^2_0(D)),\qquad F\in L^2(D;\R).
\]
Extend $w$ by zero. Suppose $\eta\in C_c^\infty(B_{2s}(x))$,
$0\le\eta\le1$, $\eta=1$ on $B_s(x)$, and
$|D^j\eta|\le C_js^{-j}$ for $j=1,2$. Then
\begin{equation}\label{eq:energy}
 \|D^2(\eta^2w)\|_2^2
 \le C\left(s^{-4}\int_{B_{2s}(x)}|w|^2
                  +\int_{D\cap B_{2s}(x)}|Fw|\right).
\end{equation}
Also, if $v\in H^2(\R^2)$ and $|\{v\ne0\}|\le a<\infty$, then
\begin{equation}\label{eq:small-support}
 \|v\|_\infty\le Ca^{1/2}\|D^2v\|_2.
\end{equation}
The constants depend only on the indicated cutoff bounds.
\end{lemma}
\begin{proof}
Multiplication by $\eta^4$ preserves $H^2_0(D)$, by approximation.
Set $v=\eta^2w$. Expansion yields
\begin{align*}
 |\Delta v|^2-\Delta w\Delta(\eta^4w)
 ={}&-8\eta^2|\nabla\eta|^2w\Delta w\\
 &+\bigl(4\eta\nabla\eta\cdot\nabla w
        +(2|\nabla\eta|^2+2\eta\Delta\eta)w\bigr)^2.
\end{align*}
Integrating the first term by parts and using
$|\nabla(\eta^2|\nabla\eta|^2)|\le C\eta s^{-3}$ gives
\begin{equation}\label{eq:intermediate-energy}
 \|\Delta v\|_2^2\le\int_D|F|\eta^4|w|
       +Cs^{-2}\int\eta^2|\nabla w|^2
       +Cs^{-4}\int_{B_{2s}(x)}|w|^2.
\end{equation}
All integrations are performed on the zero extensions. Approximation
from $C_c^\infty(D)$ in $H^2$ justifies the identities; no boundary
integral is discarded. Writing $G=\int\eta^2|\nabla w|^2$, a second integration
by parts and substitution for $\eta^2\Delta w$ give
\[
 G=-\int w\Delta v+2\int\eta w\nabla\eta\cdot\nabla w
   +\int(2|\nabla\eta|^2+2\eta\Delta\eta)w^2.
\]
Hence, for $W=\|w\|_{L^2(B_{2s}(x))}$,
$G\le W\|\Delta v\|_2+G/2+Cs^{-2}W^2$.
Substituting this bound into \eqref{eq:intermediate-energy} and applying
Young's inequality allows a multiple of $\|\Delta v\|_2^2$ to be
absorbed into the left-hand side.
The Fourier identity $\|D^2v\|_2^2=\|\Delta v\|_2^2$ proves
\eqref{eq:energy}.

For \eqref{eq:small-support}, split the Fourier integral at $R>0$.
Cauchy--Schwarz gives
\[
 \|v\|_\infty\le C(R\|v\|_2+R^{-1}\|D^2v\|_2),
 \qquad \int_{|\xi|>R}|\xi|^{-4}\,d\xi=\pi R^{-2}.
\]
Optimizing gives $\|v\|_\infty^2\le C\|v\|_2\|D^2v\|_2$.
Using $\|v\|_2\le a^{1/2}\|v\|_\infty$ and cancelling the nonzero
factor gives the claim for nonzero $v$. If the Hessian norm vanishes,
the $L^2$ function $v$ is zero. Approximation extends the Fourier
estimates to $H^2$ functions; the argument requires neither an endpoint
embedding nor a logarithmic loss.
\end{proof}

We combine the small-support inequality with the local energy bound
to estimate the clamped inverse on a sparse set.

\begin{lemma}[Weighted sparse inverse]\label{lem:inverse}
Let $\mu,A>0$ be continuous on $\Om$ and suppose
$B_{4/\mu(x)}(x)\Subset\Om$. Assume, for a fixed $K\ge1$,
\begin{equation}\label{eq:abstract-comparison}
 K^{-1}\le\frac{\mu(y)}{\mu(x)}\le K,\qquad
 K^{-1}\le\frac{A(y)}{A(x)}\le K
 \quad\text{if }|y-x|<4/\mu(x).
\end{equation}
There are $\nu_0(K)>0$ and $C(K)>0$ with the following property.
If an open set $D\subset\Om$ satisfies
\begin{equation}\label{eq:sparsity}
 |D\cap B_{2/\mu(x)}(x)|\le\nu|B_{2/\mu(x)}|\quad(x\in\Om),
 \qquad0<\nu\le\nu_0(K),
\end{equation}
and a real measurable source obeys $|F|\le H\mu^4A$ almost
everywhere on $D$ for some $H\ge0$, then there is a real
$w\in H^2_{\loc}(\Om)\cap C(\Om)$, zero on $\Om\setminus D$, such that
\begin{equation}\label{eq:inverse}
 \Delta^2w=F\quad\text{in }\mathcal D'(D),\qquad
 |w|\le C(K)\nu^2HA\quad\text{on }\Om.
\end{equation}
For finitely many admissible weights and source bounds, the same
correction can satisfy all the corresponding estimates, after taking
$\nu$ below their thresholds.
\end{lemma}
\begin{proof}
First let $D'\subset D$ be open with $\overline{D'}\Subset\Om$.
The source is bounded on $D'$ by its weighted bound. Solve
\[
 \int_{D'}\Delta w\Delta\phi=\int_{D'}F\phi,
 \qquad w,\phi\in H^2_0(D').
\]
Zero extension, the Hessian identity, and Poincar\'e's inequality on
a ball containing $D'$ imply coercivity on $H^2_0(D')$.
Lax--Milgram therefore gives a solution. By the planar $H^2$ embedding,
its zero extension is continuous and has compact support in $\Om$.
Thus $S=\sup_\Om|w|/A$ is finite and, if positive, attained at some
$x$. Set $s=\mu(x)^{-1}$. On $B_{2s}(x)$,
\[
 |w|\le KSA(x),\qquad |F|\le K^5Hs^{-4}A(x)\quad\text{on }D'.
\]
The integrals involving $w$ have support of measure at most $C\nu s^2$.
Lemma~\ref{lem:energy} gives
\[
 \|D^2(\eta^2w)\|_2^2
 \le C(K)\nu s^{-2}A(x)^2(S^2+HS).
\]
Applying \eqref{eq:small-support} to $\eta^2w$ then gives
\[
 S^2A(x)^2\le C(K)\nu^2A(x)^2(S^2+HS).
\]
Choose $\nu_0$ so that $C(K)\nu_0^2\le1/2$. Absorbing the quadratic
term and cancelling the remaining positive factor gives
$S\le2C(K)\nu^2H$. The estimate also holds when $S=0$.

To remove the compactness restriction, take
$D_j=D\cap\{j^{-1}<|x|<j\}$. The corresponding variational solutions
satisfy the pointwise bound uniformly. Equation~\eqref{eq:energy}
on fixed compact subsets of $\Om$, together with the pointwise bound,
gives uniform local $H^2$ bounds. A diagonal subsequence converges weakly
in $H^2_{\loc}$ and locally uniformly, by compact embedding into
$C^{0,\beta}$ on smaller balls for $0<\beta<1$.
The limit vanishes on $\Om\setminus D$ and inherits the pointwise
bound. Every compactly supported test function in $D$ is admissible
in $D_j$ for large $j$, so weak convergence proves the equation.
Every approximant satisfies all the finitely many weighted estimates,
so the same limit inherits them simultaneously. We do not assert the
global identity $\Delta^2w=F\mathbf1_D$ for the zero extension.
\end{proof}

\subsection{Averaging with a vanishing second moment}

We smooth the correction across the boundary of the small-value set
by an average that reproduces biharmonic functions.
Choose a real radial $\psi\in C_c^\infty(B_1^2)$ such that
\begin{equation}\label{eq:moments}
 \int\psi=1,\qquad \int|z|^2\psi(z)\,dz=0.
\end{equation}
For example, if $\rho$ is a nonnegative radial smooth unit-mass kernel,
take $\psi=(4\rho_{1/2}-\rho)/3$, where
$\rho_b(z)=b^{-2}\rho(z/b)$. This kernel is signed.
For a positive smooth radius $a(x)$, put
\[
 \mathcal S_ah(x)=\int\psi(z)h(x+a(x)z)\,dz.
\]

\begin{lemma}[Biharmonic averaging]\label{lem:smoothing}
Assume the sampling balls lie in the domain of the functions and
$|D^ja|\le C_ja^{1-j}$ for $1\le j\le4$. For locally bounded $h$,
\begin{equation}\label{eq:smoothing-jets}
 |D^j\mathcal S_ah(x)|\le Ca(x)^{-j}
                    \sup_{B_{a(x)}(x)}|h|,\qquad0\le j\le4.
\end{equation}
For smooth $f$,
\begin{equation}\label{eq:smoothing-source}
 |\Delta^2\mathcal S_af(x)|
 \le C\sup_{B_{a(x)}(x)}|\Delta^2f|.
\end{equation}
If $h$ is distributionally biharmonic on the sampling ball, then
$\mathcal S_ah(x)=h(x)$ for its smooth representative.
The constants depend only on $\psi$ and the radius derivative bounds.
Bounds of the same type at all orders give smoothness of $\mathcal S_ah$.
\end{lemma}
\begin{proof}
Write the averaging kernel as
$G(x,y)=a(x)^{-2}\psi((y-x)/a(x))$.
The chain rule gives $\|D_x^jG(x,\cdot)\|_1\le Ca(x)^{-j}$.
On each compact subset, the kernels and their derivatives have a common
compact support within the domain. Differentiation under the integral
is therefore justified for locally bounded $h$.

For a smooth biharmonic $h$, its circular mean $M_h(r)$ satisfies
\[
 M_h''+r^{-1}M_h'=\Delta h(x),\qquad
 M_h(0)=h(x),\quad M_h'(0)=0,
\]
since $\Delta h$ is harmonic. Thus
$M_h(r)=h(x)+r^2\Delta h(x)/4$ exactly. The moment conditions \eqref{eq:moments} therefore give exact
reproduction. In the distributional case, interior regularity for
constant-coefficient elliptic operators supplies the smooth
representative.

To prove \eqref{eq:smoothing-source}, let
\[
 \Phi(y)=\frac1{8\pi}|y|^2\log|y|,\quad
 T_b=\Phi*\psi_b-\Phi,\quad \psi_b(y)=b^{-2}\psi(y/b).
\]
Then $\Delta^2\Phi=\delta_0$. Choose $0<\kappa<1$ with
$\supp\psi\subset B_\kappa$. Reproduction away from the singularity
gives $\supp T_b\subset\overline{B_{\kappa b}}\subset B_b$. Also
\[
 \Phi(bY)=b^2\Phi(Y)+\frac{b^2\log b}{8\pi}|Y|^2.
\]
The additional quadratic term cancels: radiality and
\eqref{eq:moments} give
\[
 \int|Y-z|^2\psi(z)\,dz-|Y|^2
 =-2Y\cdot\int z\psi(z)\,dz+\int|z|^2\psi(z)\,dz=0.
\]
Consequently $T_b(y)=b^2T_1(y/b)$. For $q\ge1$,
$\partial_b^qT_b$ is smooth since the singular term $-\Phi$ does
not depend on $b$. Scaling gives
\begin{equation}\label{eq:kernel-bound}
 \|\partial_b^qD_y^\beta T_b\|_1
 \le C_{q,\beta}b^{4-q-|\beta|},\qquad
 q\ge1,\quad q+|\beta|\le4.
\end{equation}

Set $F=\Delta^2f$. Compact support and distributional integration by
parts give the local identity
\[
 \mathcal S_af(x)-f(x)=\int T_{a(x)}(x-y)F(y)\,dy.
\]
Differentiate with $y$ fixed, so that no derivative falls on $F$.
When $\Delta_x^2$ is applied, the terms in which the radius is not
differentiated combine \emph{distributionally} to
$(\psi_{a(x)}*F)(x)-F(x)$. Thus no $L^1$ bound is needed for the
individual fourth derivatives of the kernel.
Every other term contains
$\partial_b^qD_y^\beta T_b$, $q\ge1$, and a product
$\prod_{h=1}^qD^{\gamma_h}a$ with
$|\gamma_h|\ge1$ and $\sum_h|\gamma_h|=4-|\beta|$.
The product is bounded by $Ca^{q+|\beta|-4}$, exactly cancelling
the power in \eqref{eq:kernel-bound}. Adding $\Delta^2f=F$ proves
the estimate. Since all kernels are supported in the sampling ball,
the argument is local and requires no extension beyond the domain.
\end{proof}

\section{Completion of the real construction}\label{sec:completion}

The same variational correction will satisfy two weighted bounds:
one controls its size relative to the small-value threshold; the other
controls the smoothed residual and preserves the leading term.

\subsection{Parameter choices and the two-weight correction}

Fix all cutoff profiles, the averaging kernel, and the uniform constants
in Propositions~\ref{prop:prepared} and \ref{prop:density}.
Choose $\nu>0$ below the sparse-inverse thresholds for $E_N$ and
$E_N\mu_N^{-4}$, and small enough that $C_2\nu^2\le1/8$ in
\eqref{eq:closeness} below. The constant $C_2$ depends only on those
uniform constants and $\|\psi\|_1$.
Next choose $0<\tau<1/10$ from Proposition~\ref{prop:density}.
Fix $0<\eps\le\eps_0$ with $\eps^2\le\tau$ and an even
$L\ge\max\{L_0,16,4\eps^{-2}\}$.
The smoothing parameter $\sigma$ will be fixed below in terms of these
constants, and then $N_0$ is chosen so that $N>L$ and $N^3\ge L^4$
for every even $N\ge N_0$. Thus the order is
\[
 \text{fixed profiles}\longrightarrow\nu\longrightarrow\tau
 \longrightarrow\eps\longrightarrow L\longrightarrow\sigma
 \longrightarrow N_0,
\]
and only $N$ subsequently varies.

Apply Lemma~\ref{lem:inverse} in $D_N=D_N(\tau)$ to
$F=-\Delta^2f_N$. With weight $E_N$, \eqref{eq:small-source}
gives $H\le C(\eps^2+\tau)\le C\tau$.
With weight $E_N\mu_N^{-4}$, \eqref{eq:global-source} gives $H=M_N$.
Using the same variational approximants for both weights gives a single
real correction $w_N\in H^2_{\loc}(\Om)\cap C(\Om)$, zero off $D_N$,
satisfying
\begin{align}
 \Delta^2(f_N+w_N)&=0\quad\text{in }D_N,\nonumber\\
 |w_N|&\le C\nu^2\tau E_N,                         \label{eq:correction-small}\\
 |w_N|&\le C\nu^2M_N\mu_N^{-4}E_N.                 \label{eq:correction-order}
\end{align}
The corrected function is smooth and biharmonic in $D_N$.
Only this interior equation is used; the argument does not require a
global inhomogeneous equation for the zero extension.

Set
\begin{equation}\label{eq:smoothed}
 a_N(x)=\sigma/\mu_N(x),\qquad
 u_N^\circ=\mathcal S_{a_N}(f_N+w_N).
\end{equation}
Take $0<\sigma\le1$. Equation~\eqref{eq:frequency-jets} implies
$|D^ja_N|\le C_{j,\sigma}a_N^{1-j}$ at every fixed order, uniformly
in $N$. For example $|D^j\mu_N^{-1}|\le C_j\mu_N^{j-1}$ follows
by differentiating the reciprocal. All sampling balls avoid zero,
and $u_N^\circ$ is smooth on $\Om$.

\subsection{The dichotomy and the bounded real potential}

We separate the region where averaging reproduces the biharmonic
replacement from the region where a lower bound controls the quotient.
The gradient estimate for $f_N$, local comparison, and
\eqref{eq:correction-small} give
\begin{equation}\label{eq:closeness}
 |u_N^\circ-f_N|\le C_1\sigma E_N+C_2\nu^2\tau E_N.
\end{equation}
To see this, write $\mathcal S_{a_N}f_N-f_N$ using $\int\psi=1$
and integrate the gradient along each sampling segment. Estimate
$\mathcal S_{a_N}w_N$ by $\|\psi\|_1$.
The constants $C_1,C_2$ are therefore independent of
$\sigma,\tau,\eps,L,N$ for $\sigma\le1$; this step uses no
derivatives of the averaging kernel.
Choose $\sigma$ so that $C_1\sigma\le\tau/8$. Then
\begin{equation}\label{eq:quarter-close}
 |u_N^\circ-f_N|\le\tau E_N/4.
\end{equation}
Decrease the fixed ratio $\sigma/\tau$ further so that on every
sampling ball
\[
 E_N(y)\ge\tfrac78E_N(x),\qquad
 |f_N(y)-f_N(x)|\le\tfrac18\tau E_N(x).
\]
These follow from \eqref{eq:prepared-jets} and comparison.
If $|f_N(x)|<\tau E_N(x)/2$, then throughout the ball
\[
 |f_N(y)|<\tfrac58\tau E_N(x)<\tau E_N(y),
\]
so the sampling ball lies strictly inside $D_N$. Biharmonic
reproduction gives $u_N^\circ=f_N+w_N$ on the open set
$\{|f_N|<\tau E_N/2\}$. We obtain the dichotomy
\begin{equation}\label{eq:dichotomy}
 \begin{cases}
 \Delta^2u_N^\circ=0,& |f_N|<\tau E_N/2,\\
 |u_N^\circ|\ge\tau E_N/4,& |f_N|\ge\tau E_N/2.
 \end{cases}
\end{equation}
In particular, the function is biharmonic on an open neighborhood of
every zero in $\Om$.

By Lemma~\ref{lem:smoothing} and \eqref{eq:correction-order},
\begin{align*}
 |\Delta^2u_N^\circ(x)|
 &\le C\sup_{B_{a_N(x)}(x)}|\Delta^2f_N|
       +Ca_N(x)^{-4}\sup_{B_{a_N(x)}(x)}|w_N|\\
 &\le C_{\sigma,\nu}M_NE_N(x).
\end{align*}
Define on $\Om$
\begin{equation}\label{eq:potential}
 V_N(x)=
 \begin{cases}
 -\Delta^2u_N^\circ(x)/u_N^\circ(x),&u_N^\circ(x)\ne0,\\
 0,&u_N^\circ(x)=0.
 \end{cases}
\end{equation}
On the first set in \eqref{eq:dichotomy} it is identically zero;
on the second, $|V_N|\le C_{\sigma,\nu}M_N/\tau\le CN^3$.
All auxiliary parameters are already fixed, so $C$ is independent of
$N$. The potential is real and smooth on $\Om$. Where $u_N^\circ$ is
nonzero, it is a quotient of smooth functions; near every zero, it
vanishes identically on an open neighborhood.

\subsection{The leading term and regularity at zero}

In the fixed core, \eqref{eq:core} and \eqref{eq:frequency} give
$\mu_N=N/r$ and $E_N=|c_N|r^{N+2}$. For sufficiently small $r$ the
sampling ball stays in the core. The core polynomial is biharmonic, so
$\mathcal S_{a_N}f_N=f_N$ there. The second correction estimate,
combined with differentiation of the averaging kernel, gives
\begin{equation}\label{eq:remainder}
 |D^j\mathcal S_{a_N}w_N(x)|\le C_{N,j}|x|^{N+6-j},
 \qquad0\le j\le4.
\end{equation}
The undifferentiated bound is
$C\nu^2M_NN^{-4}|c_N|r^{N+6}$, and the $j$th derivative introduces
the factor $a_N^{-j}=(N/(\sigma r))^j$. Local comparison controls
the supremum on the sampling ball. Dependence on $N$ is allowed in
\eqref{eq:remainder}.

Let $R_N=\mathcal S_{a_N}w_N$ near zero. Its derivatives through order
four tend to zero. For $|\beta|\le3$,
\[
 \frac{|D^\beta R_N(he_i)|}{|h|}
 \le C_{N,\beta}|h|^{N+5-|\beta|}\longrightarrow0.
\]
Induction on $|\beta|$, followed by continuity at order four, shows
that the extended derivatives are the classical derivatives of the
zero extension of the remainder.
Thus $u_N^\circ\in C^4(\R^2)$ and
\[
 u_N^\circ=c_N|z|^2\operatorname{Re}(z^N)+O_N(|z|^{N+6})
\]
with the derivative bounds in \eqref{eq:remainder}.
The supremum on a small ball is at most $C_Nr^{N+2}$. Evaluation at
$(r/2,0)$ gives the lower bound $c'_Nr^{N+2}$ for sufficiently small
$r$, because $c_N\ne0$. Hence $\ord_0u_N^\circ=N+2$.

Set $V_N(0)=0$. Both $u_N^\circ(0)$ and $\Delta^2u_N^\circ(0)$
vanish: the leading polynomial is biharmonic and the fourth derivatives
of the remainder tend to zero. The equation holds pointwise everywhere.
The distributional fourth derivatives of a $C^4$ function agree with
its classical derivatives, so no residual is supported at the origin.

\subsection{Fixed normalization}

At $(1/2,0)$, $f_N=E_N=1/4$, whereas for $r\ge1$,
$E_N=2^{-L-2}r^{-L}$. Hence \eqref{eq:quarter-close} gives
\[
 |u_N^\circ(1/2,0)|\ge\tfrac14(1-\tau/4),\qquad
 \sup_{|x|\ge1}|u_N^\circ(x)|\le(1+\tau/4)2^{-L-2}.
\]
The lower bound at the interior point is strictly larger than the
upper bound outside the unit ball. Since $L$ is fixed, the envelope
is bounded; moreover, $u_N^\circ$ is continuous and tends to zero at
infinity.
Thus
\[
 K_N=\max_{x\in\R^2}|u_N^\circ(x)|>0
\]
is finite and attained inside $B_1$. Define $u_N=u_N^\circ/K_N$.
The potential and the order are unchanged, \eqref{eq:normalization}
holds, and \eqref{eq:jet} follows with $d_N=c_N/K_N\ne0$.
This proves Theorem~\ref{thm:real}.

\section{The upper estimate and sharpness in all dimensions}\label{sec:upper}

We prove the normalized upper bound independently of the construction,
then use constant extension to obtain sharpness in all three classes.

\subsection{Carleman input and normalization}

For \cite{HWZ}, we use the numbering and pagination of
arXiv:1505.05228v1. Lemma~2.1, p.~3, with $m=2$, gives a positive
nondecreasing weight
$\varpi\asymp r$ on $(0,10)$ and
\begin{equation}\label{eq:Carleman}
 \lambda^6\int\varpi(|x|)^{-1-2\lambda}|h|^2
 \le C(n)\int\varpi(|x|)^{5-2\lambda}|\Delta^2h|^2,
\end{equation}
for $h\in C_c^\infty(B_{10}\setminus\{0\})$ and
$\lambda\ge\lambda_0(n)$. Lemma~2.2, pp.~3--4, applied to the
homogeneous elliptic operator $\Delta^2$ and rescaled to a ball, gives
\begin{equation}\label{eq:interior}
 \|D^ju\|_{L^\infty(B_{s/2}(y))}
 \le C(n)s^{-j}(1+Ms^4)\|u\|_{L^\infty(B_s(y))},\qquad0\le j\le3,
\end{equation}
when $\Delta^2u=-Vu$ and $\|V\|_\infty\le M$ there.
Only boundedness of the right-hand side is used. For complex $u,V$,
apply the interior estimate to the real and imaginary parts of $u$;
their right-hand sides have modulus at most $M\|u\|_\infty$.
Adding the squared componentwise Carleman estimates gives the complex
version. No derivative of $V$ is required.

We first arrange a strict gap between the two fixed weight values:
\begin{equation}\label{eq:weight-gap}
 \varpi(21/10)<\varpi(3).
\end{equation}
Otherwise, if $\varpi(21b/10)=\varpi(3b)$ for every $b\in(0,1]$,
iteration would give $\varpi(t)=\varpi(0.7^jt)$ for $0<t\le3$,
contrary to $\varpi(r)\asymp r$. For a fixed $b$ with strict inequality,
replace $\varpi(r)$ by $b^{-1}\varpi(br)$. Applying
\eqref{eq:Carleman} to $h(\cdot/b)$ shows that the same estimate
holds with the constant multiplied by $b^{-2}$, independently of
$M,\lambda$. We use the rescaled weight below.

\begin{proof}[Proof of Proposition~\ref{prop:upper}]
Fix $p\in B_1$ and $0<r<1/2$. Take an outer cutoff equal to one
on $B_3(p)$ and supported in $B_4(p)$, and an inner cutoff zero
on $B_{r/4}(p)$ and one outside $B_{r/2}(p)$. Let their product be
$\chi$. Inner derivatives of order $j$ cost $C_jr^{-j}$; outer
derivatives have fixed bounds. The support lies in $B_5\Subset B_{10}$.
For fixed $r$, approximate $\chi u$ in $H^4$ by smooth functions with
a common compact support away from $p$. The weights are bounded on
this support, so both weighted integrals converge. Consequently,
\eqref{eq:Carleman} applies to $\chi u$ with the weight centered at $p$.

Put $S=\|u\|_{L^\infty(B_r(p))}$. If $S\ge1/2$, the claim follows
by decreasing $c$. Otherwise choose $z\in B_1$ with $|u(z)|\ge3/4$.
The interior estimate gives a Lipschitz representative on $B_6$ with
Lipschitz constant $A=C(n,C_0)(1+M)\ge1$. For $y\in B_r(p)$,
$A|z-y|\ge|u(z)|-|u(y)|\ge1/4$. With $d=(16A)^{-1}$,
$B_d(z)$ is outside $B_r(p)$, inside $B_{21/10}(p)$, and satisfies
$|u|\ge1/2$ and $\chi=1$. Therefore, with $\varpi=\varpi(|x-p|)$,
\begin{equation}\label{eq:J-bound}
 J:=\int\varpi^{-1-2\lambda}|\chi u|^2
 \ge c(n,C_0)(1+M)^{-n}\varpi(21/10)^{-1-2\lambda}.
\end{equation}

Take $\lambda\ge\max\{\lambda_0(n),n+3\}$.
The potential contribution in \eqref{eq:Carleman} is at most
$CM^2\sup_{\supp\chi}\varpi^6$ times $J$, so it is absorbed for
$\lambda\ge C(n)(1+M^{1/3})$. The commutator is
\begin{align*}
 [\Delta^2,\chi]u={}&4\nabla\chi\cdot\nabla\Delta u
 +2\Delta\chi\Delta u+4\sum_{i,j}\partial_{ij}\chi\partial_{ij}u\\
 &+4\nabla\Delta\chi\cdot\nabla u+(\Delta^2\chi)u.
\end{align*}
On the inner annulus, use \eqref{eq:interior} on balls of radius
comparable to $r$ contained in $B_r(p)$. The commutator is bounded by
$Cr^{-4}(1+Mr^4)S$. Its squared weighted integral satisfies
\[
 I_{\rm in}\le C(n)^{\lambda+1}
 r^{n-3-2\lambda}(1+M)^2S^2.
\]
On the outer annulus, fixed-radius interior estimates and
$5-2\lambda<0$ give
\[
 I_{\rm out}\le C(n,C_0)(1+M)^2\varpi(3)^{5-2\lambda}.
\]
Set $\theta=\varpi(21/10)/\varpi(3)<1$. Dividing the last bound
by $\lambda^6$ times \eqref{eq:J-bound} gives at most
\[
 C(n,C_0)(1+M)^{n+2}\lambda^{-6}\theta^{2\lambda}.
\]
Since $\log(1+M)\le C M^{1/3}$ for $M\ge1$, a fixed choice
$\lambda=T(n,C_0)(1+M^{1/3})$ absorbs both the outer contribution
and the potential term. Incorporating the fixed weight values into
the exponential constants, the remaining inequality gives
\begin{align*}
 S&\ge c(n,C_0)C(n)^{-\lambda}(1+M)^{-(n+2)/2}
                  r^{\lambda+(3-n)/2}\\
  &\ge c(n,C_0)e^{-D(n,C_0)(1+M^{1/3})}r^{\lambda+1/2}.
\end{align*}
The second inequality uses $n\ge2$ and $r<1$.
For $r<1/4$, the exponential factor is at least
$r^{(D/\log4)(1+M^{1/3})}$, proving \eqref{eq:lower-bound}.
A vanishing order strictly larger than the resulting exponent would
contradict this lower bound as $r\downarrow0$.
\end{proof}

\subsection{Extension and the three classes}

We record the preservation of essential suprema and vanishing order
under constant extension.

\begin{lemma}\label{lem:extension}
Let $2\le d\le n$, $u\in H^4_{\loc}(B_{10}^d)\cap L^\infty(B_{10}^d)$,
and $V\in L^\infty(B_{10}^d)$ solve $\Delta_d^2u+Vu=0$
distributionally. For $x\in\R^d$, $y\in\R^{n-d}$ define
$U(x,y)=u(x)$ and $W(x,y)=V(x)$. Then
\[
 \Delta_n^2U+WU=0\quad\text{in }B_{10}^n,
\]
and, for $0<r\le10$,
\begin{equation}\label{eq:extension-norms}
 \|U\|_{L^\infty(B_r^n)}=\|u\|_{L^\infty(B_r^d)},\qquad
 \|W\|_{L^\infty(B_r^n)}=\|V\|_{L^\infty(B_r^d)}.
\end{equation}
In particular, $\ord_0U=\ord_0u$.
\end{lemma}
\begin{proof}
There is nothing to extend when $n=d$. Otherwise, all derivatives
in $y$ vanish. For weak derivatives, this follows from Fubini's
theorem applied to compactly supported test functions. On compact
cylinders, the derivatives in $x$ through order four belong to $L^2$.
Hence $U\in H^4_{\loc}$, and the equation holds distributionally.
Projection of $B_r^n$ gives one inequality in each norm identity.
For the reverse inequality, take $0<t<r$ and use the positive-volume
cylinder
\[
 B_t^d\times B_{\sqrt{r^2-t^2}}^{n-d}\subset B_r^n.
\]
Fubini identifies its essential supremum with that on $B_t^d$.
Letting $t\uparrow r$ proves \eqref{eq:extension-norms}, and then
the definition of order gives the last assertion.
\end{proof}

\begin{proof}[Proof of Corollary~\ref{cor:sharp}]
Extend the real pair of Theorem~\ref{thm:real} constantly in the last
$n-2$ coordinates. Lemma~\ref{lem:extension} preserves the equation,
the normalization with $C_0=1$, the potential bound, and the order
$N+2$. The extension is $C^4$, hence in $H^4_{\loc}$.

By Proposition~\ref{prop:upper}, the exponent $1/3$ is admissible in
the largest class $\mathcal A_n^{\C}$. If $\alpha<1/3$ were admissible in the
smallest class $\mathcal A_n^{\R}$, apply \eqref{eq:admissible}
to this real family with $M=\max\{1,\|V_N\|_\infty\}$ and $p=0$.
Since $\|V_N\|_\infty\le CN^3$ by Theorem~\ref{thm:real}, it would give
\[
 N+2\le C'(1+M^\alpha)\le C''(1+N^{3\alpha})
\]
for all sufficiently large even $N$, a contradiction. Since
\[
 \mathcal A_n^{\R}(C_0,M)
 \subset\mathcal A_n^{\R,\C}(C_0,M)
 \subset\mathcal A_n^{\C}(C_0,M),
\]
all three infima equal $1/3$, with the asserted endpoint bound.
No separate complex-valued example is needed.
\end{proof}

\section{Concluding remarks}\label{sec:remarks}

\begin{remark}[Real zeros and the two weights]
Restricting to real solutions and real potentials does not improve the
upper-bound exponent. The additional difficulty in the construction
is the larger nodal set of a real angular wave. The positive envelope
regularizes amplitude cancellation, and the sparse inverse removes
the residual near the remaining real zeros. The two estimates for
$w_N$ serve different purposes: the first gives a lower bound for the
smoothed function outside the smaller small-value region; the second
preserves the prescribed polynomial jet and controls the potential
after four derivatives fall on the averaging kernel.
\end{remark}

\begin{remark}[Scope of the construction]
The potentials are real and have no prescribed sign. No uniform
H\"older or derivative bound, compact support, or smoothness at the
origin is claimed. After extension to $\R^n$, smoothness is guaranteed
off the axis $\{x_1=x_2=0\}$, not necessarily off the single origin.
The zeros need not be isolated and no boundary condition is imposed.
Applying the independently proved upper estimate with the actual
norms of the constructed planar potentials also gives
$cN^3\le\|V_N\|_\infty\le CN^3$ for all sufficiently large $N$.
This consequence is not used in the construction.
\end{remark}

\begin{remark}[Local order versus decay at infinity]
Under $u_R(x)=u(x_0+Rx)$ the potential becomes
$V_R(x)=R^4V(x_0+Rx)$, so the exponent $1/3$ corresponds algebraically to $R^{4/3}$.
The finite construction here establishes sharpness for a normalized
local vanishing-order bound. It does not by itself produce a single
solution with prescribed exponential decay in every direction at
infinity, nor does it establish the optimality of a logarithmic factor
in a Landis-type estimate. These conclusions require additional global
arguments.
\end{remark}

\begin{remark}[Dimension one]
For $u^{(4)}+Vu=0$ on an interval, with $V\in L^\infty_{\loc}$ and
$u\in H^4_{\loc}$, the vector $(u,u',u'',u^{(3)})$ is locally
absolutely continuous and solves a first-order system with locally
bounded coefficients. Gronwall's inequality gives uniqueness for zero
initial data. A nontrivial solution therefore has a nonzero derivative
of order at most three at each point, so its vanishing order is at most
three. The analogous optimal power is zero in dimension one.
\end{remark}

\begin{thebibliography}{99}
\small
\setlength{\itemsep}{3pt plus 1pt}
\setlength{\parskip}{0pt}

\bibitem{Bakri12}
L. Bakri, \emph{Quantitative uniqueness for Schr\"odinger operator},
Indiana Univ. Math. J. \textbf{61} (2012), no.~4, 1565--1580.

\bibitem{Bakri}
L. Bakri, \emph{Carleman estimates for the Schr\"odinger operator.
Applications to quantitative uniqueness},
Comm. Partial Differential Equations \textbf{38} (2013), no.~1, 69--91.
Author version: \emph{Vanishing order of solutions to Schr\"odinger equation},
\href{https://arxiv.org/abs/1111.6530v2}{arXiv:1111.6530v2}.

\bibitem{BK}
J. Bourgain and C. E. Kenig,
\emph{On localization in the continuous Anderson--Bernoulli model
in higher dimension}, Invent. Math. \textbf{161} (2005), no.~2, 389--426.

\bibitem{CG}
F. Colombini and C. Grammatico,
\emph{Some remarks on strong unique continuation for the Laplace
operator and its powers},
Comm. Partial Differential Equations \textbf{24} (1999), no.~5--6,
1079--1094.

\bibitem{Davey}
B. Davey, \emph{Quantitative unique continuation for Schr\"odinger operators},
J. Funct. Anal. \textbf{279} (2020), no.~4, Paper No.~108566.

\bibitem{DZ}
B. Davey and J. Zhu,
\emph{Quantitative uniqueness of solutions to second-order elliptic
equations with singular lower order terms},
Comm. Partial Differential Equations \textbf{44} (2019), no.~11, 1217--1251.

\bibitem{DF}
H. Donnelly and C. Fefferman,
\emph{Nodal sets of eigenfunctions on Riemannian manifolds},
Invent. Math. \textbf{93} (1988), no.~1, 161--183.

\bibitem{FGJL}
S. Feng, Z. Gan, R. Jiang, and F. Lin,
\emph{Vanishing orders of Dirichlet solutions to the Schr\"odinger equation
in dimensions three and higher},
\href{https://arxiv.org/abs/2610.01615v1}{arXiv:2610.01615v1} (2026).

\bibitem{FI}
R. L. Frank and P. Ivanisvili,
\emph{Counterexamples to the Landis conjecture in dimensions three and higher},
\href{https://arxiv.org/abs/2608.00802v1}{arXiv:2608.00802v1} (2026).

\bibitem{GL86}
N. Garofalo and F.-H. Lin,
\emph{Monotonicity properties of variational integrals, $A_p$ weights
and unique continuation},
Indiana Univ. Math. J. \textbf{35} (1986), no.~2, 245--268.

\bibitem{GL87}
N. Garofalo and F.-H. Lin,
\emph{Unique continuation for elliptic operators: a geometric-variational
approach}, Comm. Pure Appl. Math. \textbf{40} (1987), no.~3, 347--366.

\bibitem{GGS}
F. Gazzola, H.-C. Grunau, and G. Sweers,
\emph{Polyharmonic boundary value problems: Positivity preserving and
nonlinear higher order elliptic equations in bounded domains},
Lecture Notes in Mathematics, vol.~1991, Springer, Berlin, 2010.

\bibitem{HWZ}
S. Huang, M. Wang, and Q. Zheng,
\emph{Quantitative uniqueness of some higher order elliptic equations},
J. Math. Anal. Appl. \textbf{444} (2016), no.~1, 326--339.
Author version: \href{https://arxiv.org/abs/1505.05228v1}{arXiv:1505.05228v1}.

\bibitem{JK}
D. Jerison and C. E. Kenig,
\emph{Unique continuation and absence of positive eigenvalues for
Schr\"odinger operators}, Ann. of Math. (2) \textbf{121} (1985), no.~3,
463--488; appendix by E. M. Stein, 489--494.

\bibitem{Kenig}
C. E. Kenig, \emph{Some recent quantitative unique continuation theorems},
S\'eminaire \'Equations aux D\'eriv\'ees Partielles, 2005--2006,
Exp. No.~XX, \'Ecole Polytechnique, Palaiseau, 2006, 10~pp.

\bibitem{KSW}
C. E. Kenig, L. Silvestre, and J.-N. Wang,
\emph{On Landis' conjecture in the plane},
Comm. Partial Differential Equations \textbf{40} (2015), no.~4, 766--789.

\bibitem{Kuk}
I. Kukavica, \emph{Quantitative uniqueness for second-order elliptic
operators}, Duke Math. J. \textbf{91} (1998), no.~2, 225--240.

\bibitem{Lin07}
C.-L. Lin, \emph{Strong unique continuation for $m$-th powers of a
Laplacian operator with singular coefficients},
Proc. Amer. Math. Soc. \textbf{135} (2007), no.~2, 569--578.

\bibitem{LNW}
C.-L. Lin, S. Nagayasu, and J.-N. Wang,
\emph{Quantitative uniqueness for the power of the Laplacian with
singular coefficients},
Ann. Sc. Norm. Super. Pisa Cl. Sci. (5) \textbf{10} (2011), no.~3, 513--529.

\bibitem{LZ}
F.-H. Lin and J. Zhu,
\emph{Upper bounds of nodal sets for eigenfunctions of eigenvalue problems},
Math. Ann. \textbf{382} (2022), no.~3--4, 1957--1984.

\bibitem{LTY}
H. Liu, L. Tian, and X. Yang,
\emph{Quantitative unique continuation property for solutions to a
bi-Laplacian equation with a potential},
J. Lond. Math. Soc. (2) \textbf{112} (2025), no.~2, Paper No.~e70265.

\bibitem{LMNN}
A. Logunov, E. Malinnikova, N. Nadirashvili, and F. Nazarov,
\emph{The Landis conjecture on exponential decay},
Invent. Math. \textbf{241} (2025), no.~2, 465--508.

\bibitem{Meshkov}
V. Z. Meshkov, \emph{On the possible rate of decay at infinity of
solutions of second order partial differential equations},
Math. USSR-Sb. \textbf{72} (1992), no.~2, 343--361.

\bibitem{TWZ}
L. Teng, Z. Wang, and J. Zhu,
\emph{Quantitative uniqueness for bi-Laplace equations with potentials},
\href{https://arxiv.org/abs/2608.08866v2}{arXiv:2608.08866v2} (2026).

\bibitem{Zhu16}
J. Zhu, \emph{Quantitative uniqueness of elliptic equations},
Amer. J. Math. \textbf{138} (2016), no.~3, 733--762.

\bibitem{Zhu}
J. Zhu, \emph{Quantitative unique continuation of solutions to higher
order elliptic equations with singular coefficients},
Calc. Var. Partial Differential Equations \textbf{57} (2018), no.~2,
Paper No.~58.

\end{thebibliography}
\end{document}
